\section{词法分析}

在这一部分中，我们需要将输入的源程序转换为Token串，便于接下来语法分析的进行。
此外，在出现错误的情况，我们需要进行记录，并对某些错误情况进行特定的错误处理。

\subsection{总体设计}

首先，我们定义以下类：
\begin{itemize}
    \item \textbf{FileLoc}：数据类，记录Token所在的文件位置（行，行内偏移）；
    \item \textbf{Token}：数据类，记录Token类型，Token内容和Token所在的行内字符位置，并提供属性有关的查询、比较等方法；
    \item \textbf{TokenStream}：存储当前已成功读入的Token，并根据后续的语法分析需求，提供取出Token的\textbf{poll()}等方法；
    \item \textbf{Lexer}：将输入程序处理为Token串的逻辑。同时记录了分析过程中发生的错误。
\end{itemize}

编译程序通过实例化\textbf{Lexer}类，调用\textbf{analyze()}方法进行词法分析。
\textbf{analyze()}方法的返回值是一个Bool类型，指示词法分析过程中是否出错。
若返回值为真，则打印Token串，否则打印错误信息。

最后，通过调用\textbf{Lexer::getTokenStream()}，获取词法分析后得到的Token串。 
需要注意的是，词法分析完成后，无论出现错误与否，均需得到Token串并\textbf{继续执行}。
是否应停机，交由语法分析部分进行处理，详见语法分析部分。

\subsection{方法实现}

此处主要对\textbf{Lexer}内的实现进行介绍。

从文件读取字符的操作，通过创建一个内部类\textbf{Reader}实现。
\textbf{Reader}是对\textbf{BufferedReader}的包装。
整个读取过程通过\textbf{peek()}和\textbf{read()}两个操作完成。
前者只看当前待读字符，而不从文件中读出，而后者需要从文件中读出。
也就是说，本实现的字符读取过程，\textbf{不使用反压字符回文件的操作}。

\textbf{analyze()}方法通过持续调用\textbf{Reader::peek()}，
判断当前字符应采取的分析子过程并调用，直至文件末尾。

各子过程会调用\textbf{Reader::read()}从文件中实际读取字符。
子过程通过调用\textbf{Lexer::addToken()}和\textbf{Lexer::addError()}方法，记录当前过程读取成功的Token,或出现的错误。

当前peek到字符，和执行的子方法对应关系如下：
\begin{itemize}
    \item 
\end{itemize}